\documentclass[10pt,a4paper]{article}
\usepackage[utf8]{inputenc}
\usepackage[T1]{fontenc}
\usepackage[spanish,es-noshorthands]{babel}
\usepackage[margin=3cm]{geometry}
\usepackage{amsmath,amssymb}
\usepackage{booktabs,tabularx,array}
\usepackage{enumitem}
\usepackage[most]{tcolorbox}
\usepackage{xcolor}
\usepackage{fancyhdr}
\usepackage{pdflscape}
\usepackage[colorlinks=true,linkcolor=blue!50!black,urlcolor=blue!50!black]{hyperref}

\begin{document}

\begin{titlepage}
    \centering
    \vfill
    {\huge \textbf{PROCESAMIENTO DEL LENGUAJE NATURAL} \par}
    \vspace{2cm}
    {\Large Práctica 2: Clasificación de texto para análisis de sentimiento \par}
    \vspace{1cm}
    {\Large Universidad Autónoma de Madrid \par}
    \vspace{1cm}
    {\large Iván Domínguez Hernández |  ivan.dominguezh@estudiante.uam.es \par}
    \vspace{.5cm}
    {\large Lucas Miranda López | lucas.mirandal@estudiante.uam.es \par}
    \vspace{1cm}
\end{titlepage}


\tableofcontents

\newpage

\begin{thebibliography}{9}
\bibitem{bengfort2018}
B. Bengfort, R. Bilbro y T. Ojeda,
\textit{Applied Text Analysis with Python: Enabling Language-Aware Data Products with Machine Learning},
1.ª ed., O'Reilly Media, Sebastopol, CA, 2018. ISBN 978-1-4919-6304-3.

\end{thebibliography}

\end{document}
